% Conclusion and Future Work chapter
\chapter{Conclusion and Future Work}
\label{chap:conclusion}

This chapter synthesizes the project's achievements, identifies key takeaways, acknowledges limitations, and proposes directions for future research and development.

\section{Key Takeaways}
\label{sec:takeaways}

\subsection{Project success}

This project successfully developed an end-to-end machine-learning system for automated legendary Pokémon detection, from dataset construction through feature engineering, model training, threshold optimization, and interactive deployment. Key accomplishments:

\begin{enumerate}
  \item \textbf{Domain-informed feature engineering:} Designed and implemented a reproducible \texttt{FeatureBuilder} class that encodes game-domain priors (ability rates, type combinations, generation effects) alongside statistical features. This approach proved highly effective: ability-based priors alone account for 41.8\% of model importance, validating the hypothesis that domain knowledge improves imbalanced classification.
  
  \item \textbf{Class imbalance management:} Successfully handled a severe 6.9\% legendary class imbalance via stratified cross-validation, class-weighted training, and threshold optimization. The optimized Random Forest achieves F1 = 0.78 (79\% recall, 86\% precision), substantially better than the default threshold (F1 = 0.69).
  
  \item \textbf{Production-ready model:} Trained and serialized a Random Forest with supporting artifacts (\texttt{feature\_builder.joblib}, \texttt{feature\_columns.joblib}, \texttt{pokemon\_legendary\_model.joblib}), enabling reproducible inference on new data.
  
  \item \textbf{Interactive deployment:} Built a Streamlit web application and deployed it on a cost-effective DigitalOcean VPS (\$4/month), demonstrating that high-quality ML solutions need not require expensive cloud infrastructure.
  
  \item \textbf{Interpretability:} Model predictions are explainable: feature importance rankings and engineered feature snapshots allow users to understand why a Pokémon is predicted to be legendary.
\end{enumerate}

\subsection{Effectiveness of engineered features}

The dominance of ability-based legendary rates (41.8\%) and experience-related features (36.0\% combined) reveals that:
\begin{itemize}
  \item \textbf{Game mechanics encode intent:} Designers use exclusive abilities and high base experience to mark legendary Pokémon, creating strong proxy signals for classification.
  \item \textbf{Raw stats are insufficient:} Despite expectations, individual stats and composite indices (base total, bulk) contribute only 6\% to model importance. Stat distributions overlap substantially between classes.
  \item \textbf{Type combinations matter but secondarily:} Type combo legendary rates (10.3\%) are informative but weaker than ability and experience signals.
\end{itemize}

This finding has implications beyond Pokémon: in imbalanced classification on domain-specific data, encoding expert knowledge as features often outweighs raw numerical patterns.

\subsection{Model robustness}

The Random Forest achieved strong and balanced performance:
\begin{itemize}
  \item AUC-ROC = 0.974 (excellent discrimination across thresholds).
  \item Precision = 0.86 (low false positive rate; user confidence in flagged predictions).
  \item Recall = 0.79 (catches most true legends; minimal missed opportunities).
  \item Generalization: The learned type, ability, and generation priors should generalize to new Pokémon within known categories.
\end{itemize}

\section{Limitations}
\label{sec:limitations}

\subsection{Dataset constraints}

\begin{itemize}
  \item \textbf{Size:} 1025 Pokémon is relatively modest; a larger dataset (e.g., 5000+ entries with synthetic or augmented data) might improve generalization.
  \item \textbf{Class imbalance:} 6.9\% legendary class is challenging but unavoidable given the game design; oversampling techniques (SMOTE, ADASYN) could help but risk overfitting.
  \item \textbf{Feature coverage:} The dataset includes only structured attributes (stats, types, abilities). Visual features (sprite/artwork) or textual features (Pokédex descriptions) are not utilized.
\end{itemize}

\subsection{Model and threshold limitations}

\begin{itemize}
  \item \textbf{Threshold stability:} The optimal threshold (0.91) was selected on the test set. A more robust approach would use nested cross-validation or a held-out validation set for threshold selection to avoid data leakage.
  \item \textbf{Feature dependence:} Model performance relies on manually engineered features. New game mechanics or Pokémon types could render learned priors obsolete, requiring retraining.
  \item \textbf{Generalization to new generations:} While the model is trained on Generations I–IX, future generations may introduce new legendary archetypes or ability/type combinations not seen during training.
\end{itemize}

\subsection{Deployment and infrastructure}

\begin{itemize}
  \item \textbf{Scalability:} The DigitalOcean VPS is cost-effective but has limited compute/memory. High-throughput inference or batch predictions may require scaling.
  \item \textbf{Monitoring:} No automated model monitoring or drift detection is in place. Performance on new data is not tracked; periodic retraining on fresh data is manual.
  \item \textbf{API and integration:} The Streamlit app is interactive but not easily integrated into other systems (e.g., game APIs, automated pipelines). A REST API would enable broader usage.
\end{itemize}

\section{Scope for Future Improvement}
\label{sec:future_work}


\begin{itemize}
  \item \textbf{Nested cross-validation for threshold:} Implement stratified k-fold nested CV to select threshold on validation folds, ensuring robustness and avoiding test set leakage.
  
  \item \textbf{REST API:} Build a FastAPI or Flask server exposing model predictions via HTTP, enabling integration with external applications (game bots, data pipelines).
  
  \item \textbf{Multi-modal features:} Incorporate visual features (sprite images via CNN embeddings) or textual features (Pokédex descriptions via NLP embeddings) to enrich the feature set beyond structured attributes.
  
  \item \textbf{Ensemble stacking:} Combine Logistic Regression, Naive Bayes, and Random Forest via a meta-learner to potentially achieve higher F1 scores.
  
  \item \textbf{Advanced imbalance handling:} Experiment with SMOTE/ADASYN oversampling or focal loss to further improve minority-class performance.
  
  \item \textbf{Move and evolution features:} Extract and encode information about Pokémon move pools and evolution chains, which may carry legendary signals.
\end{itemize}

\section{Lessons Learned}
\label{sec:lessons}

\subsection{Domain knowledge is invaluable}

The overwhelming importance of ability-based priors (41.8\%) compared to raw stats reinforces a crucial ML lesson: \textbf{domain expertise embedded in features outweighs black-box model complexity}. Spending time understanding the problem domain (Pokémon game mechanics) yielded a simpler, more interpretable, and more effective model than relying solely on neural networks or high-dimensional embeddings.

\subsection{Threshold tuning is essential for imbalanced data}

Simply training a model with class weights is insufficient. The optimized threshold (0.91 vs. 0.5) improved F1 from 0.69 to 0.78—a non-trivial 13\% relative gain. For imbalanced classification, threshold optimization should be a standard step, not an afterthought.

\subsection{Reproducibility requires artifact serialization}

Serializing the \texttt{FeatureBuilder}, feature columns, and model together ensures that inference replicates training exactly. This is critical for production systems where feature transformations must be identical at training and serving time.

\subsection{Cost-effective deployment is viable}

A \$4/month DigitalOcean VPS is sufficient for hosting a Streamlit app serving modest traffic. For many projects, expensive cloud platforms are overkill; simple, reliable infrastructure with automated deployment is often preferable.

\section{Final Remarks}
\label{sec:final_remarks}

This project demonstrates that even on narrow, well-defined imbalanced classification tasks, thoughtful feature engineering and careful model tuning can yield production-ready systems. The combination of domain priors, statistical transformations, and ensemble learning proved robust and interpretable.

The Pokémon legendary detection problem, while playful in context, illustrates generalizable principles applicable to real-world challenges: fraud detection (rare events), medical diagnosis (rare diseases), cybersecurity (anomaly detection), and other domains where the target class is rare but high-impact.

The codebase and model artifacts are designed for reproducibility and extension. Future researchers can:
\begin{itemize}
  \item Extend the \texttt{FeatureBuilder} with new domain-specific features.
  \item Retrain the model on new Pokémon or generations.
  \item Integrate the model into applications (bots, games, analytics pipelines).
  \item Apply the feature engineering and imbalance-handling techniques to new domains.
\end{itemize}

We hope this work serves as a practical reference for building interpretable, deployable ML systems on imbalanced tabular data.
